\chapter*{MỞ ĐẦU}
\addcontentsline{toc}{chapter}{MỞ ĐẦU}
\fontsize{13}{15}\selectfont

\section*{1. Tính cấp thiết của đề tài}
Giai đoạn 2026--2030 được xem là thời kỳ bứt tốc của hệ sinh thái Internet vạn vật (IoT) và quá trình chuyển dịch sang Trí tuệ nhân tạo vạn vật (AIoT). Quy mô kết nối toàn cầu được dự báo đạt khoảng 21,9 tỷ thiết bị trong năm 2026 \citeweb{ref_iotinsider_2026} và có thể tăng lên khoảng 39 tỷ thiết bị kết nối vào năm 2030 \citeweb{ref_startus_iot_2030}. Sự gia tăng này không chỉ mở rộng số lượng nút cảm biến trong môi trường gia đình và công nghiệp, mà còn tạo áp lực lớn lên băng thông truyền tải, chi phí lưu trữ, độ trễ phản hồi và khả năng bảo vệ dữ liệu riêng tư. Vì vậy, kiến trúc hệ thống đang dịch chuyển từ mô hình xử lý tập trung trên đám mây sang mô hình xử lý phân tán tại thiết bị biên, nơi dữ liệu được phân tích gần nguồn phát sinh hơn \cite{ref_1_4}.

Trong chu kỳ phát triển 2026--2031, AIoT được dự báo tăng trưởng mạnh khi các mô hình AI được tích hợp trực tiếp với các điểm cuối IoT giàu cảm biến, chuyển hệ thống từ giám sát phản ứng sang trí tuệ tự trị tại biên \citeweb{ref_mordor_aiot_2031}. Đối với các hệ thống nhúng giám sát thiết bị gia dụng, trọng tâm không chỉ là bảo trì dự đoán dài hạn mà còn là phát hiện bất thường và cảnh báo thời gian thực ngay trên vi điều khiển. Các kỹ thuật TinyML cho phép thiết bị liên tục theo dõi chuỗi dữ liệu đa biến từ cảm biến âm thanh và rung động, sau đó đưa ra phản hồi tức thời mà không cần gửi toàn bộ dữ liệu thô lên máy chủ. Cách tiếp cận này phù hợp với xu hướng thiết kế AI bền vững, trong đó hiệu quả năng lượng và giảm tải hạ tầng tính toán được xem là tiêu chí quan trọng của các hệ thống AI hiện đại \citeweb{ref_wef_sustainable_ai_2026}.

Bên cạnh đó, từ năm 2027 trở đi, quá trình chuẩn bị cho mạng 6G và Massive IoT đặt ra yêu cầu mới về kết nối quy mô lớn, độ trễ thấp, khả năng vận hành tin cậy và bảo mật đầu cuối \citeweb{ref_ericsson_6g_massive_iot}, \citeweb{ref_bcg_6g_ai}. Sự xuất hiện của Trí tuệ nhân tạo tác nhân (Agentic AI) trong hạ tầng viễn thông cho thấy mạng kết nối trong giai đoạn tới sẽ ngày càng tự động hóa, tự tối ưu và phụ thuộc nhiều hơn vào các luồng dữ liệu từ thiết bị biên \citeweb{ref_covalense_agentic_ai_2026}. Đồng thời, khi số lượng thiết bị và dịch vụ tự động tăng nhanh, định danh máy móc và bảo mật thiết bị trở thành lớp phòng vệ quan trọng để bảo vệ toàn vẹn dữ liệu AIoT \citeweb{ref_deviceauthority_identity_2026}.

Trong bối cảnh đó, các thiết bị gia dụng có hệ thống truyền động cơ khí phức tạp như máy giặt là một đối tượng phù hợp để nghiên cứu AIoT tại biên. Các sự cố như mất cân bằng tải, lệch lồng giặt hay hư hỏng vòng bi thường biểu hiện qua rung động và âm thanh đặc trưng \cite{ref_1_3}. Nếu không được phát hiện và cảnh báo kịp thời, các lỗi này có thể gây hư hại phần cứng, làm giảm tuổi thọ thiết bị và tiềm ẩn nguy cơ mất an toàn cho người sử dụng \cite{ref_1_3}. Tuy nhiên, việc truyền liên tục dữ liệu âm thanh và rung động lên đám mây gây tốn băng thông, tăng độ trễ và tạo rủi ro riêng tư; trong khi dữ liệu lỗi thực tế lại khan hiếm, khiến các thuật toán học có giám sát dễ gặp vấn đề mất cân bằng dữ liệu \cite{ref_1_1}, \cite{ref_1_5}. Do đó, sự kết hợp giữa TinyML và mô hình học không giám sát trở thành hướng tiếp cận phù hợp, cho phép thiết bị học quy luật vận hành bình thường và phát hiện bất thường thông qua sai số tái tạo ngay tại nguồn dữ liệu \cite{ref_1_1, ref_1_4}.

Xuất phát từ những đòi hỏi thực tế và tiềm năng công nghệ trên, đề tài: \textit{"Thiết kế hệ thống AIoT phát hiện và cảnh báo bất thường cho máy giặt dân dụng sử dụng Autoencoder không giám sát trên thiết bị biên"} được lựa chọn nghiên cứu nhằm cung cấp một giải pháp giám sát độc lập, thời gian thực và bảo mật.

\section*{2. Ý nghĩa khoa học và thực tiễn}
Về mặt khoa học, dựa trên nền tảng lý thuyết về học sâu và các khung tối ưu hóa cho nền tảng biên, đề tài trực tiếp đề xuất kiến trúc mạng tự mã hóa đối xứng kết hợp kỹ thuật lượng tử hóa INT8 chuyên biệt để có thể hoạt động tối ưu trên dòng vi điều khiển ESP32-S3 \cite{ref_1_6, ref_1_7}. Đồng thời, nghiên cứu cũng đề xuất cơ chế phân mảnh trạng thái vận hành đa mô hình, tiến hành thực nghiệm và chứng minh hiệu quả vượt trội của cơ chế này trong việc nâng cao độ chính xác phân tách lỗi so với cấu trúc đơn mô hình \cite{ref_1_7}.

Về mặt thực tiễn, đồ án cung cấp một hệ thống kết nối đầu cuối hoàn chỉnh, bao gồm từ nút mạng AI tại biên xử lý đa lõi tại thiết bị cho đến ứng dụng di động thông minh cảnh báo tức thời cho người dùng. Giải pháp có tính linh hoạt cao, dễ dàng triển khai tích hợp trên các dòng máy giặt dân dụng hiện hành nhằm giảm thiểu rủi ro hỏng hóc lớn mà không cần can thiệp sâu vào cấu trúc phần cứng của nhà sản xuất.

\section*{3. Đối tượng và phạm vi nghiên cứu}
Đối tượng nghiên cứu của đồ án tập trung vào bài toán phát hiện bất thường trên chuỗi thời gian đa phương thức thông qua việc ứng dụng các kiến trúc mạng nơ-ron học không giám sát và các phương pháp xử lý tín hiệu số trong nền tảng công nghệ TinyML. 

Phạm vi nghiên cứu được xác định rõ trên ba khía cạnh. Về nền tảng phần cứng, hệ thống sử dụng vi điều khiển lõi kép Seeed Studio XIAO ESP32-S3 làm thiết bị biên, thực hiện thu thập dữ liệu qua cảm biến âm thanh kỹ thuật số INMP441 và cảm biến gia tốc ADXL345. Về nền tảng phần mềm, đề tài ứng dụng hệ điều hành thời gian thực để quản lý đa tác vụ, sử dụng giao thức MQTT cho hạ tầng truyền thông và phát triển ứng dụng di động giám sát trên nền tảng Flutter. Về môi trường thực nghiệm, các đánh giá được giới hạn trên các pha vận hành đặc trưng của máy giặt cửa trước dân dụng bao gồm các chế độ giặt nhẹ (GENTLE), giặt mạnh (STRONG) và vắt (SPIN).

\section*{4. Phương pháp nghiên cứu}
Đề tài áp dụng kết hợp hai nhóm phương pháp chính là nghiên cứu lý thuyết và kiểm thử thực nghiệm. Phương pháp nghiên cứu lý thuyết được thực hiện thông qua việc tổng hợp và phân tích hệ thống các tài liệu học thuật liên quan đến phân tích chuỗi thời gian, lý thuyết nén mô hình học sâu, công nghệ AIoT và các tiêu chuẩn truyền thông bảo mật hiện đại \cite{ref_1_1, ref_1_5, ref_1_6}.

Trong khi đó, phương pháp thực nghiệm đi theo hướng tiếp cận lấy dữ liệu làm trung tâm \cite{ref_1_8}. Quá trình này bao gồm việc tự xây dựng trạm thu thập để lấy dữ liệu thực tế tại hiện trường, đồng thời kết hợp các kỹ thuật tăng cường dữ liệu vật lý và giả lập lỗi để đánh giá giới hạn chịu đựng của mô hình. Hệ thống sau đó được triển khai và kiểm thử trên phần cứng thực tế bằng cách lập trình đa tác vụ song song trên vi điều khiển nhằm tối ưu hóa hiệu năng tính toán. Các chỉ số đánh giá học máy (AUC, F1-Score) và thông số sử dụng tài nguyên phần cứng (RAM/Flash, độ trễ) được tiến hành đo lường và thu thập trực tiếp từ các phiên chạy thực địa dài hạn \cite{ref_1_8}.

\section*{5. Cấu trúc của đồ án}
Để giải quyết các mục tiêu và nhiệm vụ đã đặt ra, ngoài phần Mở đầu, Kết luận và Tài liệu tham khảo, nội dung chính của đồ án được tổ chức thành 4 chương như sau:
\begin{itemize}
    \item \textbf{Chương 1: Tổng quan lý thuyết.} Trình bày cơ sở lý luận về bài toán phát hiện bất thường, mạng nơ-ron tự mã hóa và các giới hạn kỹ thuật.
    \item \textbf{Chương 2: Thiết kế hệ thống.} Mô tả toàn diện kiến trúc phần cứng, thiết kế phần mềm nhúng đa nhân và đề xuất lý thuyết về kiến trúc mạng tự mã hóa ba trạng thái.
    \item \textbf{Chương 3: Xây dựng dữ liệu và Huấn luyện mô hình.} Trình bày chi tiết quy trình xây dựng tập dữ liệu thực tế, kỹ thuật trích xuất đặc trưng và quá trình lượng tử hóa mô hình học máy.
    \item \textbf{Chương 4: Triển khai và Đánh giá kết quả.} Trình bày các kết quả thực thi mô hình trên vi điều khiển, đánh giá khả năng phát hiện lỗi thực tế, đồng thời thảo luận và phân tích các mặt hạn chế của hệ thống.
\end{itemize}


% ============================================================
% GHI CHÚ TÀI LIỆU THAM KHẢO CHO PHẦN MỞ ĐẦU
% ============================================================
%
% \bibitem{ref_1_1}
% P. P. Ray, "A review on TinyML: State-of-the-art and prospects,"
% \textit{Journal of King Saud University – Computer and Information Sciences}, vol. 34, pp. 1595–1623, 2022.
%
% \bibitem{ref_1_2}
% "Distributed Intelligence in the Artificial Intelligence of Things: A Comprehensive Review of Architectures, Applications, and Challenges,"
% \textit{Preprints.org}, 2026.
%
% \bibitem{ref_1_3}
% H. Mohammadi, "EARLY DETECTION OF IMBALANCE IN LOAD AND MACHINE IN FRONT LOAD WASHING MACHINES BY MONITORING DRUM MOVEMENT,"
% \textit{Master's Thesis, Sabancı University}, 2020.
%
% \bibitem{ref_1_4}
% X. Wang et al., "Empowering Edge Intelligence: A Comprehensive Survey on On-Device AI Models,"
% \textit{arXiv preprint arXiv:2503.06027v1}, 2025.
%
% \bibitem{ref_1_5}
% "Unsupervised Anomaly Detection for IoT-Based Multivariate Time Series: Existing Solutions, Performance Analysis and Future Directions,"
% \textit{PMC}, 2023.
%
% \bibitem{ref_1_6}
% "DEEP LEARNING FOR ANOMALY DETECTION: A SURVEY,"
% \textit{(Tài liệu tổng quan về các kiến trúc Deep Learning cho phát hiện bất thường)}.
%
% \bibitem{ref_1_7}
% E. Njor, M. A. Hasanpour, J. Madsen, and X. Fafoutis, "A Holistic Review of the TinyML Stack for Predictive Maintenance,"
% \textit{IEEE Access}, vol. 12, pp. 184861-184882, 2024.
%
% \bibitem{ref_1_8}
% M. Abdallah et al., "Anomaly Detection and Inter-Sensor Transfer Learning on Smart Manufacturing Datasets,"
% \textit{Sensors}, vol. 23, 2023.
%
% \bibitem{ref_startus_iot_2030}
% StartUs Insights, "IoT Report 2026: 39 Billion Connected Devices by 2030."
%
% \bibitem{ref_mordor_aiot_2031}
% Mordor Intelligence, "Artificial Intelligence of Things (AIoT) Market Size and Research Report 2031."
%
% \bibitem{ref_ericsson_6g_massive_iot}
% Ericsson, "Rethinking massive IoT for the 6G era."
%
% \bibitem{ref_bcg_6g_ai}
% Boston Consulting Group, "6G: The Network for the Future of AI and Immersive Connectivity."
%
% \bibitem{ref_covalense_agentic_ai_2026}
% Covalense Digital, "Top Six Telecom Trends 2026: Agentic AI, 6G, Open RAN, and More."
%
% \bibitem{ref_deviceauthority_identity_2026}
% Device Authority, "The State of IoT Identity Security in 2026: Why Machine Identity Is the New Perimeter."
%
% ============================================================
